%-----------------------------------------------------------------------
\section{Introduction}
\label{sec:intro}
%-----------------------------------------------------------------------

A caustic is the envelope of a family of refracted or reflected rays. Where the
envelope touches the receiving surface, neighbouring rays cross, the
surface-to-receiver map loses rank, and the ray density that a renderer
integrates becomes infinite. Every production caustic method in use today
integrates that density. The methods differ in how they estimate the integral
--- photon maps store the ray density as points \cite{jensen1996}, SPPM does so
with a radius that shrinks as it converges \cite{hachisuka2008}, path tracers
connect vertices \cite{veach1997} --- but they share the integrand, and the
integrand is singular on the very feature the user is trying to see.

This paper starts from that observation. The
first direction is theoretical, and it is the reason the paper opens with
kinematics rather than with rendering. The singularity is not an accident of
any parameterisation. It is the fold of a smooth map, it has a normal form, and
the normal form dictates both how the geometric estimate fails and what the
correct description near the fold looks like. It also fixes a number: the fold
is described near its degeneracy by a cubic phase, and the cubic phase carries
one wavelength and one geometric coefficient, whose ratio is a length. That
length scales as $\lambda^{2/3}$, not as $\lambda$ and not as the Fresnel
$\sqrt{\lambda}$. Nobody fits that exponent; it falls out of the normal form,
and it is what makes a per-channel fringe width more than a decoration.

The second direction is practical. If the wave description is right near the
fold, and the geometric description is right away from it, then a single
estimator can use each where it is valid, at the cost of detecting where the
boundary is. We implement that estimator and then spend the second half of the
paper trying to falsify it against the tools that practitioners actually run.
The evaluation is deliberately unflattering in its construction: the ground
truth is an independent brute-force trace that shares no code with the method
under test, the reference conforming solver is PBRT-v4 SPPM \cite{pharr2023},
and a path-traced comparison is produced by Blender Cycles on the same scene
with the same measured dispersive indices. Where our method loses, we say so.

\subsection{Why the kinematics come first}
\label{sec:why-kinematics}

It is possible to write a caustic renderer without ever writing down the map
that is being inverted, and most of the literature discourages doing so: photon
mapping is presented as a Monte Carlo estimator over paths, and its convergence
is discussed in terms of sample counts and radii. That presentation is
efficient, but it hides the one property that decides whether the method can
work at all. A density estimator whose target integrand diverges is not a
convergence-rate problem. Increasing the sample count does not help, reducing
the kernel radius does not help, and adding a denoiser hides the failure behind
a smooth blur that no longer means anything. The only useful question is
whether the estimator's support contains the singular set, and that question is
answered by the rank of a Jacobian.

So we write the kinematics down. Section~\ref{sec:kinematics} defines the
surface, the ray map, and the exact irradiance it induces, which is the
push-forward of the surface measure. The fold locus is then the zero set of the
Jacobian determinant, and we prove the two statements the rest of the paper
leans on: that a geometric quadrature cannot converge near the fold at any
resolution (Lemma~\ref{lem:diverge}), and that the correct local description is
the canonical fold integral whose amplitude envelope is the Airy function
(Proposition~\ref{prop:foldnormal}). From the same normal form we extract the
wavelength scaling (Corollary~\ref{cor:scaling}) and, because we have been
bitten by this, the statement that the resulting intensity is not integrable -
so a correction built on it cannot be normalised by a constant
(Lemma~\ref{lem:notl1}).

\subsection{How computer graphics builds caustics}
\label{sec:why-cg}

Section~\ref{sec:cg} then reads the standard constructions as answers to the
kinematic problem, which makes their shared limitation visible. Forward
constructions --- photon mapping and its progressive variant --- estimate the
push-forward by depositing samples and integrating density. Backward
constructions --- path tracing, and the bidirectional and vertex-connection
variants \cite{veach1997,muller2017} --- estimate it by connecting vertices
and dividing by a geometric term. Both factor the same $1/|\detJ|$, and both
become invalid exactly on $\fold$. The section states where each family is
exact, where each is merely expensive, and where neither can be repaired,
because that boundary is what the hybrid estimator is built around.

\subsection{The current way, and the way we take}
\label{sec:why-current}

Sections~\ref{sec:current} and \ref{sec:ourway} are the two halves of the
comparison. The first measures what current practice produces and costs. We
take the position that a comparative claim is only meaningful if the
comparison runs on one scene, with one set of measured material data, one
source, and one set of estimators applied identically to all sides; the section
therefore reports both the numbers and the estimator convention that produced
them, including one previously published statistic that we withdraw because it
was an artefact of that convention
(Section~\ref{sec:verify-metric}).

The second half states our method, which is small: keep the geometric
quadrature, detect the fold band analytically from $|\detJ|$, and inside the
band replace the divergent density with a normalised Airy profile while
conserving the band's own measured flux. The method introduces three
quantities --- a detector percentile, a reference fringe width, and a ceiling
--- and Section~\ref{sec:results-ablation} measures the sensitivity to each
rather than asserting a value.

\subsection{Contributions}
\label{sec:contributions}

\begin{itemize}
\item We derive the fold's kinematics from the ray map, connecting the rendering
      estimator to the fold's canonical normal form and establishing the
      $\lambda^{2/3}$ fringe scaling in the splatting pipeline (Section~\ref{sec:kinematics}).
\item The Airy fold intensity tail is non-integrable on $\mathbb{R}$, which forces
      the flux-conserving correction to be formulated as a band-local discrete
      flux ratio rather than an unphysical global constant (Lemma~\ref{lem:notl1},
      Section~\ref{sec:ours-normalisation}).
\item A hybrid estimator follows, examined through an equal-flux Gaussian ablation
      control (confirming that coarse correlation is driven by band flux bounding,
      while Airy physics controls chromatic fringe structure) and a parameter sensitivity
      study (Sections~\ref{sec:ourway}, \ref{sec:results-ablation}).
\item We evaluate like-for-like against PBRT-v4 SPPM and Blender Cycles on measured
      dispersive materials, with qualitative spectral verification against Mitsuba~3,
      and documenting estimator conventions and limitations
      (Sections~\ref{sec:results}, \ref{sec:verify}).
\item We analyze the colour of the caustic and its fringes under absorbing media
      and non-solar source spectra. Bulk absorption and emission spectra modulate
      channel throughput linearly, while the $\lambda^{2/3}$ scaling displaces fringe
      peaks chromatically across the lit band (Section~\ref{sec:results-colour}).
\item We test the next singularity stratum. The cusp ($A_3$) case is implemented
      and verified: its Pearcey diffraction integral regularizes the $|y|^{-1/3}$
      divergence, but numerical evaluation is four orders of magnitude more expensive
      than the 1D Airy lookup (Section~\ref{sec:verify-cusp}).
\item Methodological discrepancies from earlier drafts (fitted Cauchy vs Sellmeier
      dispersion, tile metric conventions, and non-circular energy audits) are
      documented with reproducible diagnostic scripts (Sections~\ref{sec:verify-index},
      \ref{sec:verify-metric}).
\end{itemize}

\subsection{What this paper does not claim}
\label{sec:notclaimed}

The correction applies to folds. A cusp requires a Pearcey function in place of
the Airy one; we build that case and measure it
(Section~\ref{sec:verify-cusp}), and the construction carries over while the cost
does not --- which is a narrower claim than we would like and a more useful one
than leaving it open. Swallowtails and the umbilic normal forms are untested, and
we make no claim about them (Section~\ref{sec:limitations}). We also do not
validate against a physical
apparatus. The comparison in this paper is entirely computational, against
solvers and against an independent trace, and the quantities compared are
therefore properties of those solvers and that scene. A laboratory measurement
would test a different claim, and the scene here is specified precisely enough
that someone may make it.

\subsection{Reading order}
\label{sec:reading}

Section~\ref{sec:kinematics} is self-contained and can be read as a short
treatment of fold optics for graphics. Readers who want only the method should
read Section~\ref{sec:ourway} after Section~\ref{sec:cg}; readers who question
the numbers should read Section~\ref{sec:estimators} first, because three of our
earlier errors lived there and not in the physics.
